\documentclass[11pt]{article}

\usepackage[utf8]{inputenc}
\usepackage[T1]{fontenc}
\usepackage{lmodern}
\usepackage[margin=1in]{geometry}
\usepackage{amsmath,amssymb,amsthm}
\usepackage{booktabs}
\usepackage{microtype}
\usepackage[hidelinks]{hyperref}
\usepackage[capitalise]{cleveref}

% Run-in headings italic, never bold.
\makeatletter
\renewcommand\paragraph{\@startsection{paragraph}{4}{\z@}%
  {3.25ex \@plus1ex \@minus.2ex}{-1em}%
  {\normalfont\normalsize\itshape}}
\makeatother

\title{The Thurstone Tilt Is a Schur Bridge,\\
and Both of Its Ends Are Races}
\author{Peter Cotton\thanks{Microprediction.
Email: \texttt{peter.cotton@microprediction.com}. A companion verification
script is \texttt{paper/verify\_tilt\_bridge.py} at
\texttt{github.com/microprediction/schur}. The identities below are checked
there in exact rational arithmetic.}}
\date{September 20, 2026}

\newcommand{\R}{\mathbb{R}}
\newcommand{\ones}{\mathbf{1}}
\newcommand{\E}{\mathbb{E}}
\DeclareMathOperator{\diag}{diag}

\newtheorem{proposition}{Proposition}
\newtheorem{corollary}{Corollary}
\newtheorem{remark}{Remark}

\begin{document}
\maketitle

\begin{abstract}
Thurstone portfolio polishing calibrates latent abilities that reproduce a
benchmark under a reference dependence law and re-races them under a richer
one, with a scalar $\phi$ dialling the estimated correlation in. Choosing the
reference to be the block-diagonal part of the estimate makes that blend
exactly the estimate with its cross-cluster entries scaled by $\phi$, and for a
two-block partition the conditional covariance of that matrix is the damped
Schur complement at $\gamma=\phi^2$. The tilt dial is therefore the Schur dial
under a square root. The pair of endpoints is new. Every bridge
catalogued so far ends at $\Sigma^{-1}u$, a linear solve that can short any
asset without bound; this one ends at a race, which lies on the simplex, splits
a duplicated position rather than doubling it, and stays bounded. The taper is
positive semidefinite for every $\phi$, so unlike an unmatched pricing space the
path has no pole, and the closed-form reliability of the cross-block estimate
transfers as $\phi^\star=\sqrt{\gamma^\star}$, which sets the dial from data
rather than by tuning.
\end{abstract}

\section{Introduction}

The bridges assembled on this site all end in the same place. Hierarchical risk
parity, hierarchical equal risk contribution, nested clustered optimization and
plain inverse variance are near ends, and turning a dial carries each of them to
$\Sigma^{-1}u$: minimum variance for $u=\ones$, maximum diversification for
$u=\sigma$, the tangency direction for $u=\mu$ \cite{cotton2024schur,
cotton2026nco, cotton2026herc}. That far end is a linear solve. It can short an
asset without bound, it doubles a position when an asset is duplicated, and it
needs a matrix inverse.

Thurstone portfolio polishing \cite{cotton2026thurstone} produces a portfolio of
a different kind. Weights are read as the winning probabilities of a race among
the assets, abilities are calibrated so that the race under a reference law
reproduces a benchmark, and the same abilities are then re-raced under a richer
law. The result is the gradient of an expected maximum, so it lies on the
simplex by construction, it splits rather than doubles a duplicated position,
and no covariance is inverted. In the Gaussian case the two laws are two
correlation matrices and a scalar $\phi\in[0,1]$ dials between them,
\begin{equation}
  C_{\mathrm{tilt}}=(1-\phi)\,C_{\mathrm{calib}}+\phi\,\widehat C ,
  \label{eq:blend}
\end{equation}
with $\phi=0$ reproducing the benchmark and $\phi=1$ using the full estimate.

This note observes that \eqref{eq:blend} is a Schur damping for the right
reference, so the tilt is a bridge in the sense of the other papers, and that
its far end is a race rather than a linear solve.

\section{The reference that makes the blend a taper}

Fix a partition of the assets into clusters and write $\ell(i)$ for the cluster
of asset $i$. The \emph{block taper} of a correlation $\widehat C$ at strength
$\phi$ leaves the diagonal blocks alone and scales everything else:
\begin{equation}
  \big(\widehat C\circ T_\phi\big)_{ij}
  =\begin{cases}\widehat C_{ij}, & \ell(i)=\ell(j),\\[2pt]
                \phi\,\widehat C_{ij}, & \text{otherwise.}\end{cases}
  \label{eq:taper}
\end{equation}

\begin{proposition}[The blend is a taper]
\label{prop:taper}
Let $C_{\mathrm{calib}}$ be the block-diagonal part of $\widehat C$, that is
$\widehat C_{ij}$ when $\ell(i)=\ell(j)$ and zero otherwise. Then for every
$\phi\in[0,1]$ the blend \eqref{eq:blend} equals the block taper
\eqref{eq:taper}.
\end{proposition}

\begin{proof}
Within a diagonal block both matrices carry $\widehat C_{ij}$, since
$(1-\phi)\widehat C_{ij}+\phi\widehat C_{ij}=\widehat C_{ij}$. Off the diagonal
blocks $C_{\mathrm{calib}}$ vanishes and the blend leaves $\phi\widehat C_{ij}$.
\end{proof}

The choice of reference matters here. The two references named in
\cite{cotton2026thurstone} are the identity and a one-factor market correlation.
Under the identity the blend damps the within-cluster correlations as well, and
is not \eqref{eq:taper}; the certificate exhibits the difference. The
block-diagonal reference is a third choice, and it is the one that aims the
damping at the cross-cluster block alone.

\section{The taper is the dial}

\begin{proposition}[Two blocks]
\label{prop:two}
Let the partition have two clusters $A$ and $B$ and let $\widehat C\succ0$. The
conditional covariance of $A$ given $B$ under the tapered correlation is the
damped Schur complement of \cite{cotton2024schur} at $\gamma=\phi^2$:
\begin{equation}
  \big(\widehat C\circ T_\phi\big)_{AA}
  -\big(\widehat C\circ T_\phi\big)_{AB}
   \big(\widehat C\circ T_\phi\big)_{BB}^{-1}
   \big(\widehat C\circ T_\phi\big)_{BA}
  \;=\;\widehat C_{AA}-\phi^2\,\widehat C_{AB}\widehat C_{BB}^{-1}\widehat C_{BA}.
  \label{eq:identity}
\end{equation}
\end{proposition}

\begin{proof}
With two clusters the taper leaves $\widehat C_{AA}$ and $\widehat C_{BB}$
unchanged and multiplies $\widehat C_{AB}$ by $\phi$. Substituting gives
$\widehat C_{AA}-\phi\widehat C_{AB}\widehat C_{BB}^{-1}\phi\widehat C_{BA}$.
\end{proof}

This is the identity of the localization note \cite{cotton2026loc} read in the
other direction. There a block-constant taper of an ensemble covariance was
shown to have a damped complement as its conditional; here the same fact says
that the tilt's blend and the Schur dial are one dial, related by
$\gamma=\phi^2$.

\begin{remark}[Beyond two blocks]
With three or more clusters the taper also damps the cross-block entries
\emph{inside} the complement $B$, so $\big(\widehat C\circ T_\phi\big)_{BB}$ is
no longer $\widehat C_{BB}$ and \eqref{eq:identity} becomes an approximation.
On random correlations of twelve assets in three clusters the two sides differ
by about $2\times10^{-5}$ at $\phi=\tfrac12$, against exact equality for two
clusters. The
flat construction of \cite{cotton2026nco}, which conditions each cluster on the
others without damping inside them, is the exact multi-cluster analogue.
\end{remark}

\section{Both ends are races}

At $\phi=0$ every cross-cluster correlation is zero, the clusters are
independent, and the race decouples into one race per cluster with the
benchmark reproduced inside each. At $\phi=1$ the taper is the estimate itself
and the race is run under the full correlation. Both ends are values of
$W_{S}(\theta)=\nabla G_S(\theta)$ with $G_S(\theta)=\E\max_i(\theta_i+\eta_i)$
\cite{cotton2026thurstone}, so both are winning-probability vectors.

That is the substance of the title. A bridge whose far end is $\Sigma^{-1}u$
inherits the far end's pathologies at every $\gamma$ near one: unbounded shorts,
a doubled position under duplication, and an inverse that must exist. A bridge
whose far end is a race inherits the opposite. The whole path lies in the
simplex, a duplicated asset splits its weight at every $\phi$, and nothing is
inverted. The dial still governs how much of the cross-cluster estimate is
believed; what changes is what belief is fed into.

\begin{corollary}[No pole]
\label{cor:nopole}
$T_\phi$ is positive semidefinite for every $\phi\in[0,1]$, so
$\widehat C\circ T_\phi$ is positive semidefinite by the Schur product theorem
and the race is well posed along the whole dial.
\end{corollary}

\begin{proof}
$T_\phi=\phi\,J+(1-\phi)\,\mathrm{blockdiag}(J)$ with $J$ the all-ones matrix,
a nonnegative combination of two positive semidefinite matrices.
\end{proof}

The contrast is worth stating, because the other direction of generalization
does not enjoy it. Pricing a cluster against something other than the portfolio
held in it makes the outer $k\times k$ matrix asymmetric, and it can pass
through singularity inside the dial \cite{cotton2026pricing}. Damping the law
the race is run under cannot: the taper is a Schur product, and positive
semidefiniteness is closed under it.

\section{Setting the dial}

Because $\gamma=\phi^2$, the closed-form reliability of the cross-block estimate
carries over. Writing $\rho$ for the conditional correlation across the block
and $n$ for the number of observations, \cite{cotton2026twosides} gives
\begin{equation}
  \gamma^\star=\frac{(n-2)\rho^2}{(n-2)\rho^2+(1-\rho^2)},
  \qquad\text{hence}\qquad
  \phi^\star=\sqrt{\gamma^\star}.
  \label{eq:phistar}
\end{equation}
The tilt strength is then read off the data rather than tuned. The square root
matters: a coupling one would damp to $\gamma^\star=\tfrac14$ on the covariance
side corresponds to $\phi^\star=\tfrac12$ on the tilt side, so a tilt that looks
mild is a heavier damping than it appears.

\paragraph{What the evidence says.} On block-structured markets with the
covariance estimated from returns, the block taper and the global tilt of the
identity reference behave differently: at two hundred and fifty observations the
block taper is the better dial, and at sixty observations neither beats racing
the clusters independently. This is the same pattern the covariance-side
experiments show, and the reason is the one \eqref{eq:phistar} encodes. The
cross-cluster correlations are the part of the estimate worth least, so a dial
aimed at them is the dial worth having.

\section{Closing}

Thurstone polishing was presented as a tilt by a view on dependence, and the
Schur papers were presented as a bridge between a heuristic and an optimizer.
With the block-diagonal reference they are one construction. The tilt's $\phi$
is the Schur $\gamma$ under a square root, its near end is the independent
cluster race, and its far end is the full race.

The far end is what makes this worth recording. Every other bridge here reaches
a linear solve, and the interest lies in stopping short of it. This one reaches
a portfolio that was already wanted for its own sake, so both ends are places
one might stand, and the dial chooses between two defensible answers rather than
between a heuristic and an optimum.

Two questions remain. The multi-cluster identity is exact only under the flat
conditioning of \cite{cotton2026nco}, and whether the cheaper taper's error
matters in practice is untested. And \eqref{eq:phistar} is derived for a
quadratic loss; whether it is the right strength when the consumer is a race,
whose loss is the conjugate of an expected maximum, is open.

\begin{thebibliography}{99}

\bibitem{cotton2024schur}
P. Cotton.
Schur complementary allocation: a unification of hierarchical risk parity and
minimum variance portfolios.
arXiv:2411.05807, 2024.

\bibitem{cotton2026twosides}
P. Cotton.
Two sides of Schur damping: high-dimensional pseudo-likelihoods and portfolio
allocation.
arXiv:2606.14798, 2026.

\bibitem{cotton2026thurstone}
P. Cotton.
Thurstone portfolio polishing: tail-sensitive Black--Litterman and beyond.
Working paper, 2026.

\bibitem{cotton2026nco}
P. Cotton.
Nested clustered optimization is one end of a Schur bridge, and the interior is
sometimes provably better.
SSRN 7480738, 2026.

\bibitem{cotton2026herc}
P. Cotton.
Hierarchical equal risk contribution is a corner of a Schur square, and two
dials lead to minimum variance.
Working paper, 2026.

\bibitem{cotton2026loc}
P. Cotton.
Covariance localization is Schur damping, and its optimal strength is a
sampling-error correction.
Working paper, 2026.

\bibitem{cotton2026pricing}
P. Cotton.
The far end of a Schur bridge does not care how its clusters are priced, and the
interior does.
Working paper, 2026.

\end{thebibliography}

\end{document}
